%!TEX program = xelatex
\documentclass[12pt,a4paper]{ctexart}

\usepackage{geometry}
\geometry{left=2.5cm,right=2.5cm,top=2.5cm,bottom=2.5cm}
\usepackage{amsmath,amssymb,amsthm}
\usepackage{booktabs,tabularx,multirow}
\usepackage{float,enumitem,xcolor}
\usepackage{hyperref}

\theoremstyle{definition}
\newtheorem{definition}{定义}[section]
\newtheorem{theorem}{定理}[section]
\newtheorem{lemma}{引理}[section]
\newtheorem{proposition}{命题}[section]
\newtheorem{corollary}{推论}[section]

\newcommand{\R}{\mathbb{R}}
\newcommand{\Z}{\mathbb{Z}}
\newcommand{\set}[1]{\{#1\}}
\newcommand{\abs}[1]{|#1|}
\newcommand{\VEC}[1]{\boldsymbol{#1}}

\title{\textbf{方案一：枚举框架 + 贪婪前向模拟\\（含停泊-再作业策略）}}
\author{数学建模备赛}
\date{2026年7月19日}

\begin{document}
\maketitle

% ==================== 问题分析 ====================
\section{问题分析}

在 $10\times 10$ 的网格海域中，船舶从起始岛 $B(1,5)$ 出发，须在 $30$ 天内抵达终点交付岛 $E(10,5)$。海域中有 $5$ 个特殊功能点：



\textbf{最优路径}（含停泊策略）：$B\to W_1\to S_2\to W_3\to S_2\to E$。

\begin{itemize}
  \item $W_1$ 作业 3 天（$+60$ Z），$W_3$ 作业 $3+1_{\text{停泊}}+3=7$ 个日历天/6 个作业天（$+168$ Z）
  \item 旅行 19 天，总天数 29 天（有 1 天裕量）
  \item $Z=100+60+168=328$
\end{itemize}

\subsection{停泊策略的收益分解}

以 $W_3$ 为例量化停泊-再作业的增量价值：
\begin{itemize}
  \item \textbf{无停泊}：每访问次最多作业 3 天，需两次造访 $W_3$ 来累计 6 天作业。中间必须往返 $S_2$（$S_2\!\to\!W_3$ 3 天，$W_3\!\to\!S_2$ 3 天），总旅行 23 天，仅剩 7 天用于作业。
  \item \textbf{含停泊}：一次造访即可作业 6 天（$3+1_{\text{停泊}}+3$），节省了 5 天往返旅行。这 5 天中的 2 天分配给 $W_1$（从 1 天增至 3 天），净增 $+40$ Z。
\end{itemize}

\subsection{DP 方法失效分析}

DP 的核心假设——到达同一状态的不同路径可合并——在本题不成立。两个不同路径可能在同一时刻到达同一位置，但持有不同的 $(O,H,F,M)$ 资源向量。一个路径可能 $Z$ 更高但 $M$ 枯竭，无法完成后续。因此"最优子结构"需同时考虑 $Z$ 和 $(O,H,F,M)$ 五维状态，状态空间爆炸导致 DP 退化为实际上不可行。

\subsection{方法评价}

\textbf{优势}：
\begin{itemize}
  \item 数学完备性——定理 \ref{thm:completeness} 提供了严格的全局最优性保证。
  \item 零依赖——纯 MATLAB 脚本，无需 Optimization Toolbox 或其他第三方库。
  \item 输出完整——包含最优路径、逐日明细（位置、动作、资源存量、采购量），满足赛题输出要求。
  \item 可解释性——三层嵌套结构清晰，每层决策可独立追踪。
\end{itemize}

\textbf{局限}：
\begin{itemize}
  \item 含停泊策略时枚举规模增大（每骨架组合从 $\sim 20$ 增至 $\sim 200$），耗时从分钟级升至数分钟。
  \item 在 $30\times 30$ 网格的任务 3 中，骨架空间将急剧膨胀（中间点更多），全枚举策略不可直接扩展。
\end{itemize}

\end{document}
